\setcounter{section}{2}

\Needspace{5\baselineskip}
\hypertarget{ring-attention}{%
\section{Ring Attention}\label{ring-attention}}

Ring Attention is a blockwise attention method for distributed multi-GPU/TPU systems. By rotating key/value blocks around a ring of devices, it distributes attention computation for extremely long sequences across devices. Published papers support this algorithmic principle, but without explicit disclosure in a model's technical report, one should not further claim that a particular commercial model uses this implementation.

FlashAttention addresses how to fit extremely long sequences into limited SRAM within a single GPU through tiling; Ring Attention addresses how to fit extremely long sequences (for example, more than 1 million tokens) into a cluster's GPU memory through block rotation across multiple GPUs, breaking through the physical memory limit of a single GPU.

Suppose you want to train a model with a context length of 10 million. Attention requires Q to interact with all K and V. Even with FlashAttention, a single H100 cannot hold such a long KV cache and the intermediate activations. Sequence parallelism is then needed to divide the long sequence across N GPUs.

Ring Attention works as follows: each GPU stores one query block (several rows of the Q matrix, that is, the queries of several tokens), while K and V blocks (each containing several matrix rows) flow among GPUs over a high-speed interconnect like ``conveyor-belt sushi.'' Each GPU processes only the small block passing through it and sends it to the next GPU when finished, never accumulating the data.

\Needspace{5\baselineskip}
This process runs for \(P\) rounds among the \(P\) nodes in the ring:

\begin{enumerate}
\def\labelenumi{\arabic{enumi}.}
\tightlist
\item
\Needspace{5\baselineskip}
  \textbf{Computational Inner Loop}:

  \begin{itemize}
  \tightlist
  \item
    GPU \(i\) uses its local \(Q_i\) and the key/value block it currently holds (denoted \(K_{\mathrm{curr}}, V_{\mathrm{curr}}\)) to calculate attention scores and local outputs.
  \item
    Computation uses FlashAttention's blockwise logic, maintaining local Softmax normalization factors and log-sum-exp statistics.
  \end{itemize}
\item
\Needspace{5\baselineskip}
  \textbf{Communication}:

  \begin{itemize}
  \tightlist
  \item
    While computing, GPU \(i\) sends its current \(K_{\mathrm{curr}}, V_{\mathrm{curr}}\) to the next GPU (GPU \(i+1\)).
  \item
    GPU \(i\) simultaneously receives a new key/value block from the previous GPU (GPU \(i-1\)).
  \end{itemize}
\item
\Needspace{5\baselineskip}
  \textbf{Overlap of Computation and Communication}:

  \begin{itemize}
  \tightlist
  \item
    The essence of Ring Attention is that matrix multiplication usually takes longer than transmitting key/value blocks, so the two can proceed in parallel.
  \item
    While the system computes block \(t\), the network is already transmitting block \(t+1\).
  \item
    If computation takes longer than transmission, communication latency is completely hidden, resulting in nearly zero additional communication overhead overall.
  \end{itemize}
\end{enumerate}

\Needspace{5\baselineskip}
The computation depends on the blockwise computation property of Softmax:

\Needspace{5\baselineskip}
Ring Attention works because Softmax can be computed blockwise, which is also the foundation of FlashAttention. Standard Softmax requires global normalization:

\[
\mathrm{softmax}(x)_i = \frac{e^{x_i}}{\sum_{j=1}^{N} e^{x_j}}
\]

However, we can split the sequence into two blocks \(A\) and \(B\), calculate their local statistics separately, and then combine them into global statistics.

\Needspace{5\baselineskip}
First calculate block \(A\)'s local maximum and local exponential sum:

\[
m_A = \max_{j \in A} x_j,\qquad
l_A = \sum_{j \in A} e^{x_j - m_A}
\]

\Needspace{5\baselineskip}
Then calculate block \(B\)'s local maximum and local exponential sum:

\[
m_B = \max_{j \in B} x_j,\qquad
l_B = \sum_{j \in B} e^{x_j - m_B}
\]

\Needspace{5\baselineskip}
Finally, combine them using simple mathematical transformations to update the global maximum \(m_{\mathrm{global}}\) and global normalization factor \(l_{\mathrm{global}}\):

\[
m_{\mathrm{global}} = \max(m_A, m_B)
\]

\[
l_{\mathrm{global}} =
e^{m_A - m_{\mathrm{global}}} l_A
+ e^{m_B - m_{\mathrm{global}}} l_B
\]

Therefore, Ring Attention only needs one \(K, V\) block in GPU memory at a time. Once the computation finishes, it updates the statistics and moves that block away.

\Needspace{5\baselineskip}
The detailed computational procedure is as follows:

For GPU \(i\), the current round uses only its local query block \(Q_i\) and the \(K_{\mathrm{curr}}, V_{\mathrm{curr}}\) it currently holds. All the following statistics are calculated separately for each query row; in matrix form, the corresponding maxima and sums are broadcast row-wise.

\Needspace{5\baselineskip}
Calculate local scores:

\[
S = \frac{Q_i K_{\mathrm{curr}}^T}{\sqrt{d}}
\]

\Needspace{5\baselineskip}
Calculate the row maxima of the current block:

\[
m_{\mathrm{block}} = \max(S)
\]

\Needspace{5\baselineskip}
Calculate the local exponential sums of the current block:

\[
l_{\mathrm{block}} = \sum \exp(S - m_{\mathrm{block}})
\]

\Needspace{5\baselineskip}
Calculate the unnormalized attention output:

\[
\tilde O_{\mathrm{block}} = \exp(S - m_{\mathrm{block}}) V_{\mathrm{curr}}
\]

The output O of GPU i is a matrix with the same number of rows as Q\_i (if each GPU stores only one query vector, O is a vector). For each query, it is the attention-weighted result over V based on the K and V blocks that have passed through that GPU so far. After each round, update the statistics and output O:

Suppose the previous round maintained \(m_{\mathrm{prev}}\), \(l_{\mathrm{prev}}\), and \(O_{\mathrm{prev}}\), while the current block provides \(m_{\mathrm{block}}\), \(l_{\mathrm{block}}\), and \(\tilde O_{\mathrm{block}}\). Online Softmax merges them as follows.

\Needspace{5\baselineskip}
Update the global maximum:

\[
m_{\mathrm{new}} = \max(m_{\mathrm{prev}}, m_{\mathrm{block}})
\]

\Needspace{5\baselineskip}
Update the normalization factor:

\[
l_{\mathrm{new}} =
e^{m_{\mathrm{prev}} - m_{\mathrm{new}}} l_{\mathrm{prev}}
+ e^{m_{\mathrm{block}} - m_{\mathrm{new}}} l_{\mathrm{block}}
\]

\Needspace{5\baselineskip}
Update the output \(O\):

\[
O_{\mathrm{new}} =
\frac{
l_{\mathrm{prev}} e^{m_{\mathrm{prev}} - m_{\mathrm{new}}} O_{\mathrm{prev}}
+ e^{m_{\mathrm{block}} - m_{\mathrm{new}}} \tilde O_{\mathrm{block}}
}{
l_{\mathrm{new}}
}
\]

In this way, each GPU does not need to store the complete \(K, V\). By merging local statistics and local outputs in every round, it obtains an attention result equivalent to global Softmax.

\FloatBarrier
\Needspace{5\baselineskip}
\hypertarget{references-93}{%
\subsection{References}\label{references-93}}

\vspace{0.25em}
\begin{itemize}
\tightlist
\item
  Liu, H., Zaharia, M., \& Abbeel, P. (2023). \href{https://arxiv.org/abs/2310.01889}{Ring Attention with Blockwise Transformers for Near-Infinite Context}. arXiv:2310.01889.
\end{itemize}

\clearpage
